\scanpage{5}
\uline{Proposition} \quad (i) Tout faisceau représentable $F\handwritten{\in} S$ est dans $\overset{\handwritten{\sim}}{S}_M$. \deleted{Si}

\deleted{on a}\inserted{(ii) Soit} un morphisme $u:F\handwritten{\to}G$ \deleted{de faisceaux,} dans $\overset{\handwritten{\sim}}{M}$, \deleted{et} Si $G\handwritten{\in}\mathrm{Ob}\,\overset{\handwritten{\sim}}{S}_M$, alors $F\handwritten{\in}\mathrm{Ob}\,\overset{\handwritten{\sim}}{S}_M$, \handwritten{et}\deleted{L}\handwritten{l}a réciproque est vraie si $u$ est couvrant.

(iii) Soit $u:F\to G$ un morphisme couvrant, tel que \deleted{$F\unsure{\in}\mathrm{Ob}\,\overset{\handwritten{\sim}}{S}_M$ et} $F\times_G F\xrightarrow{\handwritten{\mathrm{pr}_1}}F$ soit dans $\overset{\handwritten{\sim}}{M}$, alors $G$ est dans $\overset{\handwritten{\sim}}{S}_M$ sss $F$ l'est.

(i\deleted{ii}\handwritten{v}) Une somme $\handwritten{\coprod_i}F_i$ de faisdeaux est dans $\overset{\handwritten{\sim}}{S}_M$ sss les $F_i$ le sont.

(v) \deleted{Soit $I$ une catégorie finie, telle que}\inserted{Supposons} $S$ soit stable par \deleted{lim de} produit de deux objets (resp. par produit fibré de deux objets sur un \deleted{type $I$. Alors il en est de même de $\overset{\handwritten{\sim}}{S}_M$.} troisième, resp. par objet final, resp. par $\underset{\handwritten{\leftarrow}}{\lim}$ finies) alors il en est de même de $\overset{\handwritten{\sim}}{S}_M$.

(i) est trivial, \deleted{Prouvons} ainsi que la réciproque dans (ii). La partie directe est facile .... (iii) résulte de (ii), car $u$ étant couvrant, la condition $\mathrm{pr}_1\handwritten{\in} \overset{\handwritten{\sim}}{M}$ équivaut à $u\handwritten{\in} \overset{\handwritten{\sim}}{M}$. Prouvons (iv): si les $F_i$ sont dans $\overset{\handwritten{\sim}}{S}_M$, il en est de même de la somme: trivial. \deleted{\unsure{D}}Réciproque: l'inclusion $F_i\handwritten{\to} F=\handwritten{\coprod_i}F_i$ est dans $\overset{\handwritten{\sim}}{M}$ grâce à $\overset{\handwritten{\sim}}{c}$), donc si $X$ est représentable et $X\to F$ est dans $\overset{\handwritten{\sim}}{M}$, alors $X_i=X\times_F F_i \handwritten{\in} X$ est dans $\overset{\handwritten{\sim}}{M}$, donc $X_i$ peut se couvrir par des $X_{ij}\handwritten{\in} S$ avec $X_{ij}\to X_i$ dans $\overset{\handwritten{\sim}}{M}$ par $\overset{\handwritten{\sim}}{f}$); alors les composés $X_{ij}\to X_i\to F_i$ sont dans $\overset{\handwritten{\sim}}{M}$, d'autre part les $X$ envisagés couvrent $F$, donc les $X_{i\deleted{\unsure{i}}j}$ \deleted{\unsure{($X,i$)}} ($X,i$ variables) couvrent $F_i$, donc les $X_{ij}$ couvrent $F_i$, qui est donc dans $\overset{\handwritten{\sim}}{S}_M$. On laisse aussi (v\unsure{)|à} au rédacteur (attention, il est douteux que ça marche pareil pour un noyau, car il ne sera pas vrai nécessairement qu\deleted{\unsure{e}}'un \deleted{\unsure{homo}}morphisme de doublets dans $\overset{\handwritten{\sim}}{S}$ \inserted{ou même de $S$} induisant des morphismes $\handwritten{\in} \overset{\handwritten{\sim}}{M}$ sur les sommets induise un morphisme \deleted{\unsure{de}} $\handwritten{\in} M$ sur les noyaux des double flèches).

\uline{Corollaire} \quad $\mathrm{Ob}\,\overset{\handwritten{\sim}}{S}_M$ est le plus petit sous-ensemble de $\mathrm{Ob}\,S$ contenant $\mathrm{Ob}\,S$ et stable par \inserted{sommes et par} passage au quotient par une relation d'équivalence à $\mathrm{pr}_1$ qui est dans $\overset{\handwritten{\sim}}{M}$. Par suite, \deleted{\unsure{il|la}} $\overset{\handwritten{\sim}}{S}_M$ ne change pas si on remplace $S$ par une sous-catégorie $S'$ de $\overset{\handwritten{\sim}}{S}$ contenue dans $\overset{\handwritten{\sim}}{S}_M$, et $M$ par $M'=\overset{\handwritten{\sim}}{M}\handwritten{\cap}\mathrm{Fl}\,S'.$

% archive sheet number: 4
